\documentclass[11pt]{article}
\usepackage{amsmath, amssymb, graphicx, hyperref}
\usepackage{enumitem}
\setlist{nosep}
\usepackage[margin=1in]{geometry}

\title{In-Class Problem Set: Totals vs Proportions in Political and Health Data (R + GitHub)}
\author{}
\date{}

\begin{document}
\maketitle

\noindent \textbf{Goal.} Practice the difference between visualizing \emph{totals} and \emph{proportions}, and understand how these choices change interpretation. You will use state-level election and public health data to construct total counts, proportions, and grouped comparisons, then visualize them using appropriate color encodings.

\medskip
\noindent \textbf{Dataset.} The provided dataset contains one row per U.S.\ state with election results and cumulative death counts through November 7, 2020.

\medskip
\noindent \textbf{Key variables (excerpt).}
\begin{itemize}
  \item \texttt{totalvotes}, \texttt{biden\_votes}, \texttt{trump\_votes}
  \item \texttt{biden\_share}, \texttt{trump\_share}, \texttt{biden\_margin}
  \item \texttt{covid\_deaths\_to\_2020\_11\_07}
  \item \texttt{pneumonia\_deaths\_to\_2020\_11\_07}
\end{itemize}

\medskip
\noindent \textbf{What to submit (in your GitHub repo).}
\begin{itemize}
  \item A script file: \texttt{scripts/lab.R}
  \item A short write-up: \texttt{outputs/writeup.md}
  \item Saved figures in \texttt{figures/} (see requirements below)
\end{itemize}

\medskip
\noindent \textbf{Rules.}
\begin{itemize}
  \item Work inside an \textbf{R Project}.
  \item Use a \textbf{sequential, hard-coded workflow} (no user-defined functions).
  \item Save figures using \texttt{ggsave()} (no screenshots).
  \item Git commands must be run in the \textbf{Terminal tab}, not the R Console.
  \item Color choices must be intuitive and accessibility-conscious.
\end{itemize}

\section*{Questions}

\noindent \textbf{Pseudo-code expectation.} For Questions 2--6, include short pseudo-code comments in \texttt{scripts/lab.R} before each major block (for example: ``load data $\rightarrow$ derive variables $\rightarrow$ summarize $\rightarrow$ visualize").

\begin{enumerate}

  \item \textbf{Pull the data and set up your workflow (proof required).}
  \begin{enumerate}
    \item In the \textbf{Terminal tab}, run:
\begin{verbatim}
git status
git pull
\end{verbatim}
    \item Confirm the dataset exists in your repo (path specified in the course GitHub).
    \item Create the standard folder structure if missing: \texttt{scripts/}, \texttt{outputs/}, \texttt{figures/}.
    \item \textbf{Proof (write-up):} In \texttt{outputs/writeup.md}, paste:
    \begin{itemize}
      \item the output of \texttt{getwd()},
      \item the output of \texttt{list.files("data")}.
    \end{itemize}
  \end{enumerate}

  \item \textbf{Load and inspect the dataset.}
  \begin{enumerate}
    \item Load the dataset into an object named \texttt{df}.
    \item Summarize the data structure.
    \item Add 2--4 pseudo-code lines before your inspection/summarizing code.

    \textbf{Pseudo-code structure:}
\begin{verbatim}
# Step 1: load the RDS/CSV file into df
# Step 2: check dimensions and column names
# Step 3: summarize vote and death columns
\end{verbatim}

    \textbf{Suggested edit:} Use \texttt{dim(df)}, \texttt{names(df)}, and a focused summary of vote and death variables.
    \item \textbf{Proof (write-up):} Report:
    \begin{itemize}
      \item number of rows and columns,
      \item the range of \texttt{totalvotes},
      \item the range of \texttt{covid\_deaths\_to\_2020\_11\_07}.
    \end{itemize}
  \end{enumerate}

  \item \textbf{Create a total illness death variable.}
  \begin{enumerate}
    \item Create a new variable:
    \[
      \texttt{total\_illness\_deaths = covid\_deaths + pneumonia\_deaths}
    \]
    \item Compute total illness deaths aggregated across all states.
    \item Add 2--4 pseudo-code lines before coding this derivation and aggregation.

    \textbf{Pseudo-code structure:}
\begin{verbatim}
# Step 1: create total_illness_deaths column (sum of two death cols)
# Step 2: compute national total with sum()
# Step 3: find state-level min and max
\end{verbatim}
    \item \textbf{Proof (write-up):} Report the national total and the minimum/maximum state totals.
  \end{enumerate}

  \item \textbf{Visualize totals: bar plot of total illness deaths by state.}

  Create a bar plot showing \textbf{total illness deaths by state}.

  \medskip
  \noindent \textbf{Pseudo-code structure:}
\begin{verbatim}
# Step 1: reorder states by total_illness_deaths
# Step 2: build bar plot with geom_col()
# Step 3: flip coordinates for readability
# Step 4: add title, axis labels (note: totals, not proportions)
# Step 5: save with ggsave()
\end{verbatim}

  \noindent \textbf{Required:}
  \begin{itemize}
    \item Add 3--5 pseudo-code lines describing your totals bar-plot workflow.
    \item Order states by total illness deaths.
    \item Clearly label axes.
    \item State explicitly that this is a \emph{total}, not a proportion.
  \end{itemize}

  \noindent Save as:
  \[
    \texttt{figures/total\_illness\_deaths\_by\_state.png}
  \]

  \item \textbf{Visualize proportions: Trump vs Biden.}

  Create a categorical variable indicating the election winner:
  \begin{itemize}
    \item \texttt{winner = "Biden"} if \texttt{biden\_share > trump\_share}
    \item \texttt{winner = "Trump"} otherwise
  \end{itemize}

  \textbf{Pseudo-code structure:}
\begin{verbatim}
# Step 1: create winner variable based on biden_share vs trump_share
# Step 2: group by winner and compute proportion of total illness deaths
# Step 3: build bar plot with proportions (geom_col)
# Step 4: apply red/blue color scale
# Step 5: add labels clarifying these are proportions
\end{verbatim}

  Then:
  \begin{enumerate}
    \item Add 3--5 pseudo-code lines describing your winner classification and proportional plot workflow.
    \item Create a bar plot showing the \textbf{proportion of total illness deaths} accounted for by Trump- vs Biden-won states.
    \item Use an appropriate color scheme:
    \begin{itemize}
      \item red--blue for categorical winner, or
      \item red--purple--blue if you choose to encode \texttt{biden\_margin}.
    \end{itemize}
  \end{enumerate}

  \noindent \textbf{Required:}
  \begin{itemize}
    \item Colors must be intuitive and colorblind-friendly.
    \item The plot must clearly communicate that values are \emph{proportions}, not raw counts.
  \end{itemize}

  \noindent Save as:
  \[
    \texttt{figures/illness\_deaths\_by\_winner\_proportion.png}
  \]

  \item \textbf{Faceted visualization: deaths by state, grouped by winner.}

  Using \texttt{facet\_wrap()} or \texttt{facet\_grid()}, create a faceted bar plot that shows total illness deaths by state, split into separate panels by \texttt{winner} (Biden vs Trump).

  \medskip
  \noindent \textbf{Pseudo-code structure:}
\begin{verbatim}
# Step 1: reorder states by total_illness_deaths within each facet
# Step 2: build bar plot with geom_col()
# Step 3: facet by winner using facet_wrap() or facet_grid()
# Step 4: use scales = "free_y" so each panel has its own state list
# Step 5: add appropriate colors, title, axis labels
# Step 6: save with ggsave()
\end{verbatim}

  \noindent \textbf{Required:}
  \begin{itemize}
    \item Use \texttt{facet\_wrap(\textasciitilde winner, scales = "free\_y")} or an equivalent \texttt{facet\_grid()} call.
    \item States should be ordered by death count within each panel.
    \item Use color to encode winner (red/blue) consistently with Question 5.
    \item Clearly label the plot so it is interpretable on its own.
  \end{itemize}

  \noindent Save as:
  \[
    \texttt{figures/illness\_deaths\_faceted\_by\_winner.png}
  \]

  \item \textbf{Interpretation (write-up required).}

  In \texttt{outputs/writeup.md}, write 12--16 sentences addressing:
  \begin{itemize}
    \item How interpretation changes when moving from totals to proportions.
    \item What the bar plot of totals emphasizes that the proportional plot hides.
    \item What the proportional plot emphasizes that the totals plot hides.
    \item How the faceted plot compares to the single proportional plot---what does faceting reveal or obscure?
    \item Why your color choices are appropriate for the task.
    \item One way these plots could be misinterpreted if shown without context.
  \end{itemize}

  \item \textbf{Submit your work (GitHub or Canvas; proof required).}
  \begin{enumerate}
    \item Choose \textbf{one} submission path:
    \begin{itemize}
      \item \textbf{GitHub path:} In the \textbf{Terminal tab}, run:
\begin{verbatim}
git status
git add .
git commit -m "Totals vs proportions lab"
git push
\end{verbatim}
      \item \textbf{Canvas path:} Upload \texttt{scripts/lab.R}, \texttt{outputs/writeup.md}, and required files from \texttt{figures/} to Canvas.
    \end{itemize}
    \item \textbf{Proof (write-up):} Paste:
    \begin{itemize}
      \item if using GitHub: the output of \texttt{git status} after committing and \texttt{git log -1},
      \item if using Canvas: a short note confirming upload date/time and the list of uploaded files.
    \end{itemize}
  \end{enumerate}

\end{enumerate}

\section*{Optional challenge (if you finish early)}
Create an alternative visualization where you encode \texttt{biden\_margin} as a continuous color scale (red--purple--blue) instead of the binary winner variable. In 5--7 sentences, explain which version is clearer and for what audience.

\section*{Checklist (before you leave)}
\begin{itemize}
  \item \texttt{scripts/lab.R} runs top-to-bottom
  \item Required figures exist in \texttt{figures/}
  \item \texttt{outputs/writeup.md} includes interpretation + proofs
  \item Work is submitted (either pushed to GitHub or uploaded to Canvas)
\end{itemize}

\end{document}
